% Readers and new formats: the design (one reader per format, four data models), the models
% (recording index, ECG record and its JSON form, DICOM objects, catalog tables), the readers that
% exist, adding a format of each kind, and the two formats read by code outside the repository.
% Sources: docs/readers.md; checked against viewers/ecg/readers.js (its header comment),
% libs/engine/src/sources/mod.rs, medorama/src/formats/mod.rs, medorama/src/server.rs (page!).
\chapter{Readers and new formats}
\label{chap:readers}

This chapter explains how Medorama and its viewers read files and databases, and how to add a
format that they do not read yet. It is written for integrators and contributors who need to
bring in a new format, and for reviewers who want to know where format-specific code lives.
Adapters, the separate programs that read licensed recording formats, are described in
\cref{chap:addons}; the index format in \cref{chap:formats-reference}.

\screenshot{readers.pdf}{Readers: one per format, each producing one of four data
models.}{fig:readers-models}

\section{The design}
\label{sec:readers-design}

Every file format has its own \textbf{reader}, and every reader produces one of four
\textbf{data models}: a recording's index, an ECG record, a DICOM object, or rows of the
catalog's canonical tables. Past the reader, nothing depends on the format. VitalStream draws an
index the same way whether it came from a WFDB record or an EDF file; cohorts search entries the
same way whether they came from a CSV file or an OMOP database; and the privacy gate checks the
same fields whatever their origin.

This design has three consequences:
\begin{itemize}
\item \textbf{A new format is one new reader.} No viewer, cohort or privacy code changes, as long
  as the reader produces one of the four models.
\item \textbf{Readers can be written in any language.} A reader in another language plugs in
  through a documented boundary: the adapter protocol for recordings, the JSON form for ECGs,
  DICOM for images, and a table or a set of SQL views with the canonical columns for tabular
  data.
\item \textbf{Identity stays out of the models.} No data model has a place for the patient's
  identity. A recording's index, an ECG record and the viewers' output never hold a name, an
  identifier or a birth date. Identifiers reach Medorama only through the catalog's identifier
  tables (imported columns, connected views, and the \code{identify} action for identifiers
  stored inside recordings), where the privacy gate controls them.
\end{itemize}

\section{The four data models}
\label{sec:readers-models}

\begin{xltabular}{\linewidth}{@{}P{2.3cm}P{4.4cm}P{3.4cm}L@{}}
\caption{The four data models.}
\label{tab:readers-models}\\
\toprule
\tabhead Model & \tabhead What it is & \tabhead Specification & \tabhead Used by \\
\midrule
\endfirsthead
\toprule
\tabhead Model & \tabhead What it is & \tabhead Specification & \tabhead Used by \\
\midrule
\endhead
\bottomrule
\endfoot
Recording index & Channels, pieces of contiguous samples, min/max envelopes, vital-sign series,
  events, markers and breaks & \Cref{chap:formats-reference};
  \file{viewers/vitalstream/docs/index-format.md} & VitalStream; Medorama (a Holter's heart rate,
  recordings in cohorts and on patient pages) \\
ECG record & Leads in millivolts, their sampling rate, and the report (measurements, statements,
  notes); also its JSON form & The header comment of \file{viewers/ecg/readers.js}
  (\cref{sec:readers-ecg-record}) & The ECG viewer \\
DICOM object & Frames, measurements and reports, as the DICOM standard defines them & DICOM
  (PS3) & The PACS viewer; the ECG viewer for DICOM ECGs \\
Catalog tables & Patients, encounters, bed stays and entries, in canonical columns &
  \Cref{chap:importing} (the column reference) and \cref{chap:connecting} (the canonical views)
  & Patient pages, cohorts, the \ui{Databases} tab \\
\end{xltabular}

\subsection{Recordings: the index}
\label{sec:readers-index}

A recording (an ICU monitor's waveforms, a Holter, an EEG, a sleep study) is turned into an
\textbf{index}: a folder that holds the recording's structure (its channels and the pieces each
channel is made of), min/max envelopes of every waveform for fast drawing over days, and its
markers, vital-sign series, alarms and breaks. The index never holds the recording's samples in
full. When a view zooms in close, the samples are read from the recording itself, through the
same reader.

The index format is specified so that independently written programs produce the same bytes.
The three VitalStream viewers (Rust, C++ and Python) and Medorama share index folders, and a test
checks that they write byte-identical indexes. The language-neutral model behind the index
(channels, pieces, lanes, breaks and ``time not stated'') is described in
\cref{chap:formats-reference}.

\subsection{ECG records}
\label{sec:readers-ecg-record}

The ECG viewer reads every format into one \textbf{record}, defined in the header comment of
\file{viewers/ecg/readers.js}:
\begin{lstlisting}[style=text]
{ name, start, fs, leads: [{ name, units, fs, v }], meta }
\end{lstlisting}

\begin{xltabular}{\linewidth}{@{}P{2.0cm}L@{}}
\caption{The fields of an ECG record.}
\label{tab:readers-record}\\
\toprule
\tabhead Field & \tabhead Content \\
\midrule
\endfirsthead
\toprule
\tabhead Field & \tabhead Content \\
\midrule
\endhead
\bottomrule
\endfoot
\code{name} & What the record calls itself (a WFDB record's name, \code{MUSE ECG}). The page
  shows the file's name instead. \\
\code{start} & When the ECG was recorded, in Unix seconds (UTC), or \code{null} when the file
  does not say. \\
\code{fs} & The sampling rate in Hz. Each lead may state its own \code{fs}. \\
\code{leads} & The leads, in the file's order. \code{name} is the lead's label (\code{I},
  \code{II}, \code{aVR}, \code{V1}, \code{ECG II}; the twelve standard leads are found by name),
  \code{units} is \code{mV}, and \code{v} holds the samples in millivolts (\code{NaN} where a
  sample is missing). \\
\code{meta} & Optional. It may hold \code{fields} (label and text pairs), \code{measurements}
  (\code{name}, \code{value}, \code{unit}, and optionally \code{lead}), \code{statements}
  (\code{source}: \code{computer}, \code{confirmed} or \code{edited}, with optional
  \code{program}, \code{by} and \code{time}, and \code{lines}), and \code{notes}. \\
\end{xltabular}

The patient's identity is never part of a record. \code{ECG.checked} in
\file{viewers/ecg/ecg.js} refuses a record that the paper cannot draw: one with no leads, no
samples or no sampling rate.

\paragraph{The JSON form.} The same record can be written as a JSON file, which the viewer opens
like any other ECG:
\begin{lstlisting}[style=text]
{ name, start, fs, leads: [{ name, fs?, units, samples: [number | null] }], meta? }
\end{lstlisting}
Here \code{start} is Unix seconds, ISO 8601 text or \code{null}; \code{units} is \code{mV},
\code{uV} (microvolts) or \code{V}; a missing sample is \code{null}; and \code{meta} is checked,
so that unknown parts and identifying fields are dropped. A program in any language can convert a
format the page does not read into this form. \cmd{muse-py json} does so for GE MUSE XML.

\subsection{DICOM objects}
\label{sec:readers-dicom}

Images (echocardiography and other ultrasound loops, angiography, CT, MR, nuclear medicine and
X-ray), structured reports, encapsulated PDF documents and DICOM waveforms (ECGs) keep their own
standard model, DICOM. No other image model is defined. An image format that is not DICOM is
either shown as it is (pictures, videos, PDF documents) or converted to DICOM. The PACS viewer's
reader is \file{viewers/pacs/dicom.js}, with \file{jpeg.js} for lossless and 12-bit JPEG. The
privacy gate rewrites DICOM headers in \file{medorama/src/dicom.rs}.

\subsection{Catalog tables}
\label{sec:readers-catalog}

Tabular data (patients, encounters, bed stays, results, drugs, conditions, notes, and the files
they name) end up in the catalog's \textbf{canonical columns}. Imported tables reach them through
the recognized column names and mappings; known exports reach them through an import format that
converts them first; and connected databases reach them through three SQL views
(\code{persons}, \code{visits} and \code{entries}). The columns are specified in
\cref{chap:importing} and \cref{chap:connecting}.

\section{The readers that exist}
\label{sec:readers-existing}

\begin{xltabular}{\linewidth}{@{}P{3.4cm}P{3.0cm}P{1.8cm}L@{}}
\caption{The readers of Medorama and its viewers, and the model each one produces.}
\label{tab:readers-existing}\\
\toprule
\tabhead Format & \tabhead Reader & \tabhead Language & \tabhead Where, and the model \\
\midrule
\endfirsthead
\toprule
\tabhead Format & \tabhead Reader & \tabhead Language & \tabhead Where, and the model \\
\midrule
\endhead
\bottomrule
\endfoot
WFDB records, multi-segment records, collections, folders of records & The WFDB reader, through
  wfdb-x (refused in a build without WFDB support: \cref{sec:installation-no-wfdb}) & Rust, C++,
  Python & \file{libs/engine/src/sources/wfdb\_folder.rs},
  \file{viewers/vitalstream/cpp/engine/src/sources\_wfdb.cpp},
  \file{vitalstream\_py/sources/wfdb\_folder.py}. Recording index \\
EDF, EDF+, BDF, BDF+ & The EDF reader & Rust, C++, Python &
  \file{libs/engine/src/sources/edf.rs}, \file{cpp/engine/src/sources\_edf.cpp},
  \file{vitalstream\_py/sources/edf.py}. Recording index (\cref{sec:formats-reference-edf}) \\
BedMasterEx \file{.Stp} (GE Unity and Carescape, Philips PIICiX and Viridia) & stp-adapter, a
  private add-on available to institutions under agreement & Rust, C++, Python & Outside this
  repository, through the adapter protocol (\cref{chap:addons}). Recording index \\
WFDB (a \file{.hea} with its signal files) & \file{wfdb.js} & JavaScript & \file{viewers/ecg/}.
  ECG record \\
GE MUSE XML & \file{muse.js}; \cmd{muse-py} (\cmd{json}, \cmd{wfdb}, \cmd{index}) & JavaScript;
  Python & \file{viewers/ecg/}, \file{viewers/ecg/python/muse/}. ECG record (and its JSON form, a
  WFDB record, a table) \\
DICOM ECG waveforms & \file{dicom-ecg.js} with \file{dicom.js}; \cmd{dicom-py wfdb} & JavaScript;
  Python & \file{viewers/ecg/}, \file{viewers/pacs/python/dicom/}. ECG record (or a WFDB
  record) \\
EDF, EDF+, BDF, BDF+ & \file{edf.js} & JavaScript & \file{viewers/ecg/}. ECG record \\
The ECG viewer's JSON form & \code{ECG.readJson} in \file{ecg.js} & JavaScript &
  \file{viewers/ecg/}. ECG record \\
PNG, JPEG, PDF & Shown as they are & & \file{viewers/ecg/}, \file{viewers/pacs/}. Not a model \\
DICOM images, loops, series, structured reports, encapsulated PDF & \file{dicom.js},
  \file{jpeg.js}; \cmd{dicom-py index} & JavaScript; Python & \file{viewers/pacs/}. DICOM object
  (and a table) \\
DICOM headers in Privacy Mode & \file{dicom.rs} & Rust & \file{medorama/src/}. DICOM object,
  de-identified \\
CSV, Parquet, Excel, OpenDocument tables & \file{tables.rs}, \file{parquet\_rows.rs} & Rust &
  \file{medorama/src/}. Catalog tables \\
GE MUSE metadata tables; CLIF tables & \file{formats/ge\_muse.rs}; \file{formats/clif.rs} & Rust
  & \file{medorama/src/formats/}. Catalog tables \\
A site's own exports & TOML descriptions (\file{formats/declared.rs}) & None: a description & A
  folder of \code{formats\_dirs}. Catalog tables \\
OMOP CDM 5.3 and 5.4 & Built-in views (\file{sources/omop.rs}) & SQL & \file{medorama/src/sources/}.
  Catalog tables (connected) \\
Any other SQL database & Three views that you write & SQL & A TOML file. Catalog tables
  (connected) \\
\end{xltabular}

\section{Adding a format}
\label{sec:readers-add}

Choose the data model first, because it decides where the reader goes. A format of waveforms that
people scroll through for hours is a recording. A 10-second, 12-lead ECG printed on paper is an
ECG record. An image is DICOM. Anything tabular becomes catalog tables.

\subsection{Recordings: a built-in reader or an adapter}
\label{sec:readers-recordings}

A recording format can be read in two ways, compared in \cref{tab:readers-builtin-adapter}.

\begin{xltabular}{\linewidth}{@{}P{2.4cm}LL@{}}
\caption{A built-in reader or an adapter.}
\label{tab:readers-builtin-adapter}\\
\toprule
 & \tabhead Built-in reader & \tabhead Adapter \\
\midrule
\endfirsthead
\toprule
 & \tabhead Built-in reader & \tabhead Adapter \\
\midrule
\endhead
\bottomrule
\endfoot
What it is & A module of the recording engine, written three times: in Rust, C++ and Python & A
  separate program, named \file{<name>-adapter}, installed beside Medorama and the viewers \\
Language & Rust, C++ and Python, all three & Any language \\
Distributed & With Medorama and the viewers, in this repository & Separately, as an add-on,
  which may be licensed or private \\
Choose it when & The format is open and public, worth reading everywhere without an
  installation, and its rules can be written down exactly (as \cref{sec:formats-reference-edf}
  does for EDF) & The format is licensed or restricted, needs a vendor's library, or you want
  one implementation in a language of your choice \\
\end{xltabular}

\paragraph{A built-in reader.} The EDF reader, \file{libs/engine/src/sources/edf.rs}, is the
model to follow; the module table is in \file{libs/engine/src/sources/mod.rs}.
\begin{enumerate}
\item \textbf{Write down the reading rules} of the format exactly, as
  \cref{sec:formats-reference-edf} does for EDF: the header, the numbers, the start time, how
  pieces and breaks are made, the events, every key of \file{meta.json}, and the notes. The three
  implementations must agree byte for byte, so the rules must leave nothing to interpretation.
\item \textbf{Write the Rust module} in \file{libs/engine/src/sources/}. It provides a test of the
  file (\code{is\_<format>}), its format label, \code{build(path, progress, cancelled) -> (meta,
  arrays)} for the index, and a \code{Reader} whose \code{read(piece, a, b)} returns the physical
  samples \code{a} to \code{b - 1} of a piece. The module never reads the patient's identity into
  the index.
\item \textbf{Register it} in \file{sources/mod.rs}, with a line in \code{kind\_of},
  \code{format\_label} and \code{build}, and in the \code{Reader} enum (\code{new} and
  \code{read}).
\item \textbf{Write the same reader in C++ and in Python.} The C++ reader is
  \file{viewers/vitalstream/cpp/engine/src/sources\_<format>.cpp}, declared in
  \file{include/engine/sources.hpp}, dispatched in \file{sources.cpp} and listed in the engine's
  \file{CMakeLists.txt}. The Python reader is \file{vitalstream\_py/sources/<format>.py},
  dispatched in \file{sources/\_\_init\_\_.py}.
\item \textbf{Add test files}, synthetic only, to \file{viewers/vitalstream/testdata/}, written by
  its \file{generate.py}, and tests in each language (\file{libs/engine/tests/} and the Python
  viewer's tests).
\item \textbf{Check that the three agree.} \cmd{make parity} indexes the test files with every
  built viewer and checks that the indexes are byte-identical.
\item \textbf{Teach Medorama the format} where it needs to know it: \code{is\_recording} in
  \file{medorama/src/tables.rs}, so that an imported row makes a recording; the file kinds in
  \file{medorama/src/patient/files.rs} and \file{medorama/src/gate.rs} if the format's files are
  also sent to the browser (in Privacy Mode the gate refuses any kind it cannot de-identify); and
  \file{medorama/src/identify.rs} if the file stores identifiers, so that the \code{identify}
  action can read them into the identifier tables.
\end{enumerate}

\paragraph{An adapter.} The protocol is described in \cref{sec:addons-protocol}.
\begin{enumerate}
\item Write a program named \file{<name>-adapter} that answers \cmd{describe} (its formats, each
  with a kind, a label and extensions), \cmd{detect}, \cmd{index PATH DIR} (the index files, in
  the index format), \cmd{serve} (a session that returns samples), \cmd{identify} (the
  identifiers the file states) and \opt{-{}-version}.
\item Test it with every viewer. The following command indexes a recording through the adapter
  with each built viewer and checks that the results are byte-identical:
\begin{lstlisting}[style=shell]
viewers/vitalstream/tests/parity/index_parity.py --adapters-dir DIR FILE...
\end{lstlisting}
\item Install it with \cmd{medorama addons install PROGRAM} and restart Medorama. Nothing in
  Medorama or the viewers names the adapter's format: until the adapter is installed, its files
  are treated like any other file that the programs do not read.
\end{enumerate}

\subsection{ECGs: a reader in the ECG viewer, or a converter}
\label{sec:readers-ecg}

\paragraph{A reader in the ECG viewer.} The viewer keeps one registry of readers,
\file{viewers/ecg/readers.js}, and each format has its own file (\file{wfdb.js}, \file{muse.js},
\file{dicom-ecg.js}, \file{edf.js}).
\begin{enumerate}
\item \textbf{Write the reader in its own file.} It maps the format to the record of
  \cref{sec:readers-ecg-record}, and never reads the patient's identity. \file{ecg.js} provides
  the helpers: units to millivolts, metadata from ``key: value'' lines, and \code{ECG.checked}.
\item \textbf{End the file with} \code{Readers.add(\{ key, label, extensions, sniff, read \})}.
  \code{sniff(name, bytes, type)} says whether a file's content is of this format.
  \code{read(name, bytes, sibling)} returns a promise of the record, and \code{sibling(fileName)}
  gives a file next to this one, such as a WFDB header's signal files. Three fields are
  optional: \code{companions} (extensions of files read with it and not opened alone),
  \code{parts} (the files of a folder that belong to this record) and \code{loose} (its sniff is a
  guess, tried after the others).
\item \textbf{Load it} in \file{viewers/ecg/index.html}, after \file{readers.js}. The order of the
  scripts is the order in which the page lists the formats.
\item \textbf{List it} in Medorama's built-in page files (\code{page!("ecg", \dots)} in
  \file{medorama/src/server.rs}).
\item \textbf{Test it} in \file{viewers/ecg/tests/}; \file{readers.test.js} checks what the
  registry finds for a file.
\item \textbf{Privacy Mode.} For Medorama to send such files in Privacy Mode, the privacy gate
  must be able to de-identify them (\code{file\_kind} in \file{medorama/src/gate.rs}). Otherwise
  they are refused in Privacy Mode, which is the safe default.
\end{enumerate}

\paragraph{A converter in any language.} A program can convert a format the page does not read
into the ECG viewer's JSON form, which the page opens like any other ECG. The privacy gate,
however, refuses JSON files in Privacy Mode, because it cannot check what they hold. For ECGs that
will be viewed through Medorama in Privacy Mode, convert to a WFDB record instead, as
\cmd{muse-py wfdb} and \cmd{dicom-py wfdb} do: the gate de-identifies WFDB headers.

\subsection{Images: convert to DICOM}
\label{sec:readers-images}

The PACS viewer reads DICOM, and shows pictures, videos and PDF documents as they are. An image
format that is neither should be converted to DICOM with a conversion tool that keeps its
acquisition date, modality and pixel spacing. DICOM is also the safer choice for Medorama. In
Privacy Mode the gate rewrites a DICOM file's header (names, IDs, dates and UIDs), whereas
pictures and videos are sent as they are, with the ``may contain PHI'' warning, because the gate
cannot know what they show.

\subsection{Tables: an import format or a connected database}
\label{sec:readers-tables}

Tabular data need a reader only when their columns cannot simply be mapped. The options, in order
of effort, are as follows.
\begin{enumerate}
\item \textbf{A mapping}, with no reader at all. Map the columns once in the \ui{Import} tab, or
  with \opt{-{}-mapping} (\cref{chap:importing}).
\item \textbf{A TOML description}, for an export that your site imports regularly. It describes
  the export column by column, with date forms, value maps, numbers and joined columns, and needs
  no code and no build (\file{medorama/src/formats/declared.rs}; \cref{chap:importing}).
\item \textbf{A Rust converter}, for what a description cannot express, such as one row that
  becomes several entries, or values that depend on other rows or tables. It implements the two
  traits of \file{medorama/src/formats/mod.rs}:
  \begin{itemize}
  \item \code{Format}: \code{key} (its name on the command line and in the page), \code{label},
    \code{category}, \code{about}, \code{detect(headers)} (whether a table with these columns is
    of this format) and \code{start(headers)} (its converter);
  \item \code{Converter}: \code{headers()} (the converted table's columns, which are catalog
    fields or columns of its own), and \code{row(record)} or \code{rows(record)} (the converted
    row, several rows, or the reason a row is left out).
  \end{itemize}
  Register it in \code{builtin()} in the same file. \file{ge\_muse.rs} (one row per ECG) and
  \file{clif.rs} (one format per CLIF table, with rows that become several entries) are the
  examples. A format keeps only what it names. Identifiers that it maps to the catalog's patient
  columns are masked in Privacy Mode, and it should not keep other identifying columns, such as
  staff names or account numbers.
\item \textbf{A connected database's views}, when the data should stay in their database: three
  SQL queries with the canonical columns, or the built-in OMOP views (\cref{chap:connecting}). A
  new kind of database engine is added in \file{medorama/src/sources/} and
  \file{medorama/src/sql/}.
\end{enumerate}

\section{Two exceptions: WFDB and BedMasterEx .Stp}
\label{sec:readers-exceptions}

Two formats are read by code that does not live in this repository.
\begin{itemize}
\item \textbf{WFDB} is read through \textbf{wfdb-x}, a public WFDB library with Rust, C++ and
  Python implementations (\url{https://github.com/alphanumericslab/wfdb-x}). Each of the three
  engines uses the library of its own language, fetched at a pinned version (\file{Cargo.toml},
  \file{viewers/vitalstream/cpp/cmake/deps.cmake}). The engines' WFDB modules add what WFDB itself
  does not define: collections, folders of records read as one session, and the index
  (\cref{chap:recordings}).
\item \textbf{BedMasterEx \file{.Stp}} files are read by the \textbf{stp-adapter}, a private
  add-on available to institutions under agreement, which plugs in through the adapter protocol
  (\cref{chap:addons}). Medorama and the viewers contain no code for this format.
\end{itemize}
